\documentclass[11pt,a4paper]{article}

% Page geometry
\usepackage[margin=1in]{geometry}

% Font: XeTeX + CJK
\usepackage{fontspec}
\usepackage{xeCJK}
\setCJKmainfont{Songti SC}[BoldFont={Heiti SC}, ItalicFont={STKaiti}]
\setCJKsansfont{Heiti SC}
\setCJKmonofont{Heiti SC}
\setmainfont{Times New Roman}
\setmonofont{Menlo}[Scale=0.85]

% Tables
\usepackage{booktabs}
\usepackage{array}
\usepackage{tabularx}
\usepackage{longtable}

% Colors (needed by hyperref color expressions)
\usepackage{xcolor}

% Hyperlinks
\usepackage[colorlinks=true,linkcolor=blue!60!black,citecolor=blue!60!black,urlcolor=blue!60!black]{hyperref}

% Code listings
\usepackage{listings}
\lstset{
  basicstyle=\ttfamily\small,
  frame=single,
  breaklines=true,
  breakatwhitespace=false,
  columns=fullflexible,
  keepspaces=true,
  xleftmargin=1em,
  xrightmargin=1em,
  aboveskip=1em,
  belowskip=1em,
}

% Graphics
\usepackage{graphicx}

% Page numbering and headers
\usepackage{fancyhdr}
\pagestyle{fancy}
\fancyhf{}
\fancyfoot[C]{\thepage}
\renewcommand{\headrulewidth}{0pt}

% Paragraph spacing
\usepackage{parskip}

% Title formatting
\usepackage{titlesec}
\titleformat{\section}{\Large\bfseries}{}{0em}{}
\titleformat{\subsection}{\large\bfseries}{}{0em}{}

% For the abstract environment
\renewenvironment{abstract}{%
  \begin{center}
    \bfseries 摘要
  \end{center}
  \itshape
}{\par\bigskip}

% -----------------------------------------------------------------------
\begin{document}

\begin{center}
  {\LARGE\bfseries Metanet 网络：\\[0.3em]
  一个激励检索的去中心化 CDN}\\[1.5em]
  {\large Alex Tong}\\[0.3em]
  \texttt{alex@bitfs.org}\\[0.8em]
  {\small 版本 0.0.1 --- 2026 年 3 月}\\[1.5em]
\end{center}

\begin{abstract}
我们提出一种去中心化内容分发网络，其中节点通过提供数据服务赚取收入，而非通过证明存储数据。Metanet 网络是一条 BSV 同构侧链，为 BitFS 去中心化文件系统提供激励和分发层。热门内容通过自组织的市场机制在网络中复制：热门文件产生 payment 检索费，吸引更多节点缓存，从而提升可用性并吸引更多用户。冷门内容通过内容所有者与节点之间简单的一对一存档合约保存，合约由基于 ECDH 的存储证明验证——该证明在毫秒内完成，而非零知识替代方案所需的数小时 GPU 计算。双币种模型将终端用户与 CDN 的内部经济机制隔离：用户以 BSV 支付，而 MNT 代币仅在网络运营商之间流通。最终构建出一个随需求扩展、奖励有用工作、且无需中心化协调的内容分发网络。
\end{abstract}

% -----------------------------------------------------------------------
\section{1. 引言}

去中心化内容分发问题至今未被解决。现有系统各自解决了其中一些环节，但没有任何一个系统构建出经济激励与用户实际需求——快速、可靠地访问数据——完全一致的网络。

IPFS 证明了内容寻址的点对点分发在技术上是可行的，但它没有为节点提供经济激励去服务它们自身不需要的内容。数据可用性取决于利他主义或原始上传者的持续在线。当热门内容消失时，网络没有机制来恢复它。

Filecoin 通过付费让存储提供者持有数据来弥补激励缺口，但其激励结构围绕存储而非检索。提供者通过证明自己随时间持续存储数据来赚取奖励，使用计算密集的零知识证明——每个扇区需要数小时的 GPU 计算。相比之下，检索市场只是一个附带考虑，经济权重极小。无论是否有人下载过数据，存储提供者都能获得奖励。结果是一个为归档而非分发优化的网络。

Arweave 采用不同方法，收取一次性费用实现永久存储。虽然简洁优雅，但 Arweave 没有经济机制鼓励快速检索或地理分布。数据是永久的，但访问延迟没有保障。

传统 CDN（如 Cloudflare 和 Akamai）高效分发内容，但它们是中心化服务，可以审查、封禁，且存在单点故障风险。它们还要求与内容提供者建立信任关系。

我们提出 Metanet 网络——一个反转 Filecoin 模型的去中心化 CDN。我们不再付费让节点存储数据，而是让节点通过服务数据来赚取收入。核心洞察很简单：如果一个节点能从分发某个文件中获利，它就会选择缓存该文件。如果一个文件很热门，许多节点都会缓存它。网络围绕需求自组织，产生类似传统 CDN 的行为，但无需任何中心化协调。

对于不够热门、无法通过检索费维持自身的内容，所有者可以与单个节点签署存档合约，直接付费要求它们维护副本。这些合约使用基于 ECDH 重加密和 Merkle 挑战的存储证明，用毫秒级操作替代 Filecoin 的 zk-SNARK 机制。

设计哲学可归结为一句话：反转 Filecoin 模型。不要付费让节点存储数据，让节点通过服务数据来赚钱。热数据通过市场激励自组织。冷数据使用简单的一对一合约。一切运行在 Bitcoin Script 之上。

% -----------------------------------------------------------------------
\section{2. 与 BitFS 的关系}

Metanet 网络与 BitFS 是互补系统，设计上协同工作，但也能独立运行。两者的关系类似于 IPFS 与 Filecoin，但在设计和范围上都有重要区别。

BitFS 是构建在 BSV 区块链上的去中心化加密文件系统协议。它提供 Unix 风格的文件系统语义（目录、文件、符号链接和硬链接）、基于 Method~42 ECDH 密钥派生的默认加密、通过哈希时间锁合约 (HTLC) 实现的去信任数据交易，以及 Agent 优先的 HTTP 接口。BitFS 定义了文件是什么、谁拥有它、如何加密、以及如何转移访问权。BitFS 白皮书~[1] 详细描述了该协议。

Metanet 网络负责文件存储在哪里以及如何到达请求它们的用户。它是位于 BitFS 之下的分发和激励层，提供内容可用性、地理分布和节点运营商的经济激励。用户与 BitFS 交互；Metanet 网络提供使 BitFS 文件可访问的基础设施。

内容所有者使用 BitFS 发布文件，将文件元数据（所有权、加密参数、定价）记录在 BSV 区块链上。然后 Metanet 网络将加密的文件内容分发到节点网络中，节点通过微支付向请求者提供服务。所有者无需运行任何基础设施：一旦文件发布且元数据上链，Metanet 网络就能分发它。

这种分离很重要。BitFS 是任何服务器都可以实现的协议。用户可以在自己的机器上运行 BitFS daemon，直接提供文件服务，完全不需要接触 Metanet 网络。Metanet 网络面向那些希望获得去中心化、永远可用的内容分发服务而又不想维护自己服务器的用户。两个系统共享设计理念和密码学原语，生成独立的二进制文件：\texttt{bitfs} 用于文件系统协议，\texttt{metanet} 用于节点运营。BitFS 使用共享核心库 libbitfs-go；Metanet 目前独立实现，计划未来整合共享库。

% -----------------------------------------------------------------------
\section{3. 三层架构}

完整系统跨三个层运行，每层提供不同的保障并服务于不同的需求。

\begin{lstlisting}
 第一层: BSV 主链
 +--------------------------------------------------------+
 |  文件所有权 (Metanet DAG), HTLC 购买,                    |
 |  payment 下载, 元数据永久存储                                |
 |  货币: BSV                                              |
 +--------------------------------------------------------+

 第二层: 自托管 Daemon
 +--------------------------------------------------------+
 |  个人内容服务器 (bitfs daemon),                           |
 |  直接文件服务, 所有者完全控制                              |
 |  货币: 无 (所有者自有基础设施)                             |
 +--------------------------------------------------------+

 第三层: Metanet Chain
 +--------------------------------------------------------+
 |  去中心化 CDN, 存档合约,                                  |
 |  支付通道, 挖矿奖励                                      |
 |  货币: MNT 代币                                          |
 +--------------------------------------------------------+
\end{lstlisting}

第一层是 BSV 区块链，作为整个系统的信任根基。所有文件所有权记录、Metanet DAG 结构和 HTLC 购买交易都存储在这里。该层是永久的、抗审查的，由 BSV 网络的全部算力保障安全。

第二层是自托管的 BitFS daemon。任何文件所有者都可以在自己的服务器上运行 daemon，直接向请求者提供内容。这是最简单的部署模式：所有者保持完全控制，只需支付自己的托管成本。不需要任何代币。

第三层是 Metanet Chain，本文的核心主题。它是一条专用区块链，协调去中心化 CDN，管理内容所有者与 Metanet 节点之间的存档合约，并处理流媒体和微支付场景的支付通道。

用户根据需求选择存储和分发模式。小型永久文件可以直接存储在 BSV 上。所有者控制下的大文件使用自托管 daemon。需要去中心化高可用分发的文件使用 Metanet 网络。这些方式可以自由组合：所有者可以将元数据存储在 BSV 上，通过 Metanet 网络分发热门内容，同时自托管冷门存档。

两条链的职责清晰划分。BSV 主链记录文件所有权（Metanet DAG），处理买卖交易（HTLC 原子交换和 payment 支付），并永久存储元数据，面向终端用户和 AI Agent。Metanet Chain 记录存档合约、存储证明、支付通道结算和挖矿奖励，面向内容所有者和节点运营商。数据永远不需要在两条链之间跨链传输——它们在同一生态系统中扮演互补角色。

% -----------------------------------------------------------------------
\section{4. Metanet Chain 设计}

Metanet Chain 是一条 BSV 同构区块链。"同构"意味着结构完全相同：Metanet Chain 使用相同的交易格式、相同的 Script 引擎、相同的 UTXO 模型和相同的密码学原语。唯一的区别是独立的创世块、调整后的区块参数，以及新增的合并挖矿支持。

这个设计选择是刻意的。通过保持与 BSV 交易格式和 Script 引擎的完全兼容，Metanet Chain 不需要自定义虚拟机，不需要新的交易类型，也不需要新的脚本语言。Metanet Chain 上的所有合约——包括存档合约、存储证明和支付通道——都用标准 Bitcoin Script 实现。任何能构造 BSV 交易的工具只需最小修改就能构造 Metanet Chain 交易。

实际的实现路径很直接：fork 现有 BSV 节点软件，替换创世块，调整区块大小和间隔参数，添加辅助工作量证明 (AuxPoW) 支持以实现合并挖矿。其他一切保持不变。

合并挖矿允许 Metanet Chain 与 BTC 和 BSV 矿工共享 SHA-256 工作量证明。矿工将 Metanet Chain 区块哈希嵌入其 BTC 或 BSV 的 coinbase 交易中。Metanet Chain 通过验证以下内容来确认区块：父链的区块头满足 Metanet Chain 的难度目标，coinbase 交易包含在父链的 Merkle 树中，且 coinbase 包含正确的 Metanet 区块哈希。这一机制允许 Metanet Chain 从现有的 SHA-256 挖矿生态系统中借用安全性，而无需专用算力。

为防止对 Metanet Chain 的长程攻击，周期性的锚定交易被写入 BSV 区块链。大约每 100 个 Metanet Chain 区块，一个汇总该范围区块的 Merkle 根被发布在 BSV 的 \texttt{OP\_RETURN} 交易中。这些锚定创建了将 Metanet Chain 历史绑定到 BSV 安全性的检查点。试图重写 Metanet Chain 历史的攻击者还需要重写相应的 BSV 锚定，这在经济上是不可行的。

% -----------------------------------------------------------------------
\section{5. MNT 代币经济学}

Metanet Chain 拥有自己的原生货币 MNT 代币，其供应计划镜像比特币：总供应量 2100 万枚，初始区块奖励 50~MNT，每 210,000~块减半，目标区块时间约 10~分钟。

双币种模型是核心设计特征。两种货币服务于两个不同的市场：

BSV 是面向用户的货币。当用户通过 Metanet 网络下载文件时，他们以 BSV 支付 payment 检索费。当用户通过 HTLC 购买加密内容的访问权时，他们以 BSV 支付。BitFS 的普通用户永远不需要获取、持有或知道 MNT 代币的存在。

MNT 是面向运营商的货币。当内容所有者希望确保其数据无论热门与否都在网络中保持可用时，他们以 MNT 支付 Metanet 节点来持有存档合约。挖矿奖励以 MNT 计价。节点间数据批发交易使用 MNT。

这种分离存在是因为两个市场具有不同的动态。用户对内容的需求以 BSV 计价，这是他们已经持有的货币。节点运营商对 CDN 托管合约的需求是一个内部市场，受益于专用代币——其价值反映 CDN 服务本身的供需关系，而非更广泛的 BSV 经济。

Metanet 节点有三种收入来源：用户下载内容时支付的 BSV payment 检索费，内容所有者为确保数据可用性而支付的 MNT 存档合约费，以及协议为出块支付的 MNT 挖矿奖励。

MNT 的价值直接源于这一结构。希望获得数据可用性保障的内容所有者购买 MNT 来支付存档合约。因此，MNT 的需求是 CDN 托管服务需求的直接函数。随着网络增长、更多内容所有者寻求去中心化分发，MNT 的需求相应增长。

% -----------------------------------------------------------------------
\section{6. 热数据：自组织 CDN}

Metanet 网络的核心贡献在于：热门内容不需要存储合约，不需要存储证明，也不需要协议管理的复制。它完全通过市场激励自组织。

机制如下运作。当用户从 Metanet 节点下载文件时，他们支付 BSV 的 payment 检索费。这笔费用是节点的收入。监控网络的节点运营商观察到某些文件持续产生检索费用。理性的反应是缓存这些文件，因为服务它们是有利可图的。随着文件变得更热门，更多节点独立决定缓存它。更多缓存副本意味着更好的可用性、通过地理分布实现的更低延迟，以及更强的弹性。更好的可用性吸引更多用户，产生更多费用，吸引更多节点。

\begin{lstlisting}
热度上升  -->  payment 收入增加  -->  更多节点缓存该文件
    ^                                          |
    |                                          v
更多用户使用  <--  可用性提升  <--  网络中更多副本
\end{lstlisting}

这个正反馈循环产生一个自组织的 CDN。网络对任何给定文件的复制因子完全由市场力量决定。一个病毒式传播的视频可能被全球数百个节点缓存，一篇小众学术论文可能只被少数几个节点缓存。在两种情况下，复制水平都与需求成正比，这是经济效率最优的结果。

这种方法有四个值得注意的特性。热数据不需要存储合约，因为节点自愿缓存是有利可图的。不需要存储证明，因为 BSV 区块链上的 payment 交易记录本身就证明了节点拥有数据。不需要副本管理，因为市场自动决定正确的副本数量。不需要中心化协调者，因为每个节点基于自身的成本收益分析做出独立的缓存决策。

Metanet 节点通过三种渠道之一获取缓存数据。内容所有者可以直接推送数据到网络。节点可以从所有者的 daemon 拉取数据，支付 payment 费用。或者节点可以通过 MNT 支付通道从另一个节点批发购买数据，使位于网络良好连接区域的节点能够高效地再分发热门内容。

节点的决策逻辑很直接：如果某个文件的预期 payment 收入超过存储和服务成本，就缓存该文件；否则不缓存。这不需要协议强制执行。做出良好缓存决策的节点赚取更多收入，做出糟糕决策的节点浪费资源。市场选择高效行为。

这个模型还自然地处理了内容发现问题。当用户请求文件时，他们首先从 BSV 区块链（或通过 SPV 从所有者的 daemon）解析文件的 Metanet 元数据。元数据通过内容哈希标识文件。然后用户通过两种机制之一定位持有该文件的 Metanet 节点：查询 Metanet Chain 上的活跃存档合约（列出签约节点），或查询跟踪哪些节点自愿缓存文件的 Overlay 网络查找服务。用户根据延迟、价格或信誉选择最佳可用节点，支付 payment 费用下载。这个路由过程的每一步都是去中心化的。

% -----------------------------------------------------------------------
\section{7. 冷数据：存档合约}

并非所有数据都是热门的。个人备份、合规性存档或很少被访问的数据集不会产生足够的检索费来吸引自愿缓存。对于这类内容，Metanet 网络提供存档合约：内容所有者与 Metanet 节点之间直接的一对一协议。

存档合约是 Metanet Chain 上的 Bitcoin Script 交易。所有者将 MNT 代币锁定到一系列 UTXO 中，每个证明期一个。在每个期间，Metanet 节点必须提交有效的存储证明来领取该期的付款。如果节点未能在截止日期前提交证明，所有者可以取回锁定的代币。

存储证明机制建立在 BitFS 用于文件加密的相同 ECDH 原语之上。当所有者在存档合约下向 Metanet 节点分发数据时，数据使用所有者密钥与节点密钥之间的 ECDH 共享秘密派生的密钥进行重新加密。这意味着每个节点持有的是密码学上唯一的数据副本。节点~A 持有的副本使用与节点~B 不同的密钥加密。它们无法互相复制对方的数据来满足存储证明。\footnote{当前 Metanet 实现使用 HMAC-SHA256 模拟 ECDH 共享密钥派生（\texttt{shared\_secret = HMAC-SHA256(owner\_priv, node\_pub)}），因 Metanet 独立于 BSV SDK。安全属性等效于 ECDH 的对称密钥派生，但不提供前向保密性。生产版本将迁移至真实 secp256k1 ECDH。}

验证使用 Merkle 挑战-响应协议。每个证明期有一个从合约交易 ID 确定性派生的挑战：\texttt{challenge = SHA256(contract\_txid || period\_number)}。这决定了节点必须呈现哪个数据块。节点以块数据及其 Merkle 证明作为响应，任何人都可以根据合约中记录的 Merkle 根进行验证。由于合约的交易 ID 在合约创建前是不可预测的，节点无法预知哪些块会被挑战，因此必须保留所有块。

副本策略完全由所有者决定。协议不强制要求最少副本数。想要三份数据副本的所有者与三个不同节点分别签署三份存档合约。想要地理冗余的所有者选择不同区域的节点。协议提供机制；所有者决定策略。

这种方法与 Filecoin 的存储证明形成鲜明对比。Filecoin 要求复制证明 (PoRep) 和时空证明 (PoSt)，两者都基于 zk-SNARK。Filecoin 的密封过程每个 32~GB 扇区需要数小时的 GPU 计算。Metanet 网络中的 ECDH 重加密和 Merkle 证明验证在毫秒内完成。代价是 Metanet 网络的证明更简单，验证特定块的持有而非整个数据集的连续存储。对于一个主要目的是内容分发而非归档保证的 CDN 来说，这是设计空间中恰当的位置。

\begin{table}[htbp]
\centering
\small
\begin{tabularx}{\textwidth}{@{}l >{\raggedright\arraybackslash}X >{\raggedright\arraybackslash}X@{}}
\toprule
 & \textbf{Filecoin} & \textbf{Metanet 网络} \\
\midrule
副本独立性 & zk-SNARK (PoRep) & ECDH 重加密 \\
持续验证 & zk-SNARK (PoSt) & Merkle 挑战-响应 \\
计算成本 & GPU 密集，每扇区数小时 & 毫秒级 ECDH + AES \\
副本管理 & 协议强制最少副本 & 所有者决定（签 N 份合约） \\
检索激励 & 间接（检索市场存在但非主要激励） & 强（payment 为主要激励） \\
代币用途 & 存储 + 检索 + 质押 & 仅 CDN 托管 + 挖矿（用户用 BSV） \\
\bottomrule
\end{tabularx}
\caption{存储证明对比：Filecoin vs.\ Metanet 网络}
\end{table}

% -----------------------------------------------------------------------
\section{8. 内容收益分成}

Metanet 网络在内容所有者和节点运营商之间创建了一个自然的市场。当 Metanet 节点服务一个文件并收取 payment 检索费时，问题随之而来：谁保留收入？

内容所有者在文件元数据中设置 \texttt{revenue\_share} 参数，指定服务节点保留多少百分比的 payment 费用以及多少百分比回流给所有者。典型的分成可能是 70\% 给节点、30\% 给所有者，但比例完全由所有者自行决定。

这产生了所有者和节点之间的两种合作模式。对于热门内容，收益分成模式占主导。节点免费缓存和服务内容，赚取其份额的检索费。所有者无需运行任何基础设施即可从热门内容中获得被动收入。对于冷门内容，付费托管模式适用。所有者通过存档合约付费给节点以保证可用性，接受内容不会产生有意义的检索收入这一现实。

两个极端之间存在一个连续谱。所有者可以设置慷慨的收益分成（比如 90\% 给节点）来吸引对中等热门内容的缓存。拥有高度热门内容的所有者可以设置较低的分成（比如 50/50），因为知道节点仍会觉得有利可图。市场动态自然决定均衡比例。如果所有者将节点份额设得太低，没有节点会缓存内容。设得太高，所有者会不必要地放弃收入。

Metanet 节点沿以下逻辑评估每个文件：如果预期检索收入乘以节点份额超过存储和带宽成本，则在收益分成模式下缓存该文件；如果所有者提供存档合约，评估合约付款是否覆盖成本；否则拒绝存储该文件。不需要任何协调——每个节点独立进行此计算。

% -----------------------------------------------------------------------
\section{9. 支付通道}

虽然单次 payment 支付对于离散的文件下载效果良好，但流媒体和大文件传输需要更高效的支付机制。即使 BSV 费用很低，为视频流的每个千字节支付单独的链上交易也是不现实的。

Metanet 网络为这些场景支持支付通道。支付通道是一种双方安排：资金锁定在链上的 2-of-2 多签交易中，后续支付作为已签名但未广播的余额分配更新进行交换。只有最终的结算交易需要上链，将可能数千次的微支付压缩为单笔交易。

网络中运行两种类型的支付通道。BSV 支付通道连接用户与 Metanet 节点，实现流媒体和大文件下载的按千字节计费。MNT 支付通道连接内容所有者与节点（用于持续的存档合约付款）以及节点与节点之间（用于数据批发购买）。

payment HTTP 协议扩展了通道感知的头部。客户端表明对通道支付的支持，服务器响应通道定价（可能低于单次支付价格，反映减少的每次交易开销）。开放通道内的后续请求包含签名的余额更新而非离散支付。服务器验证签名，更新本地余额状态，并立即提供内容。单个请求不需要链上确认。

争议解决遵循标准支付通道惯例。每次余额更新都携带递增的序列号。双方交换先前状态的撤销密钥。如果一方广播过时的状态，另一方可以在定义的时间窗口内提交惩罚交易，声索作弊方的全部通道余额。这一机制由 \texttt{OP\_CHECKSEQUENCEVERIFY} 时间锁保障，确保诚实行为而无需信任。

支付通道支持几个重要的使用场景：按秒计费的流媒体，按千字节计费的大文件下载（中断的下载可以在不丢失通道余额的情况下恢复），以及建模为周期性通道支付的长期订阅。

% -----------------------------------------------------------------------
\section{10. 与现有系统对比}

下表总结了 Metanet 网络与现有去中心化存储和分发系统以及传统 CDN 的对比。

\begin{table}[htbp]
\centering
\small
\setlength{\tabcolsep}{4pt}
\begin{tabularx}{\textwidth}{@{}l >{\raggedright\arraybackslash}X >{\raggedright\arraybackslash}X >{\raggedright\arraybackslash}X >{\raggedright\arraybackslash}X >{\raggedright\arraybackslash}X@{}}
\toprule
 & \textbf{Metanet} & \textbf{Filecoin} & \textbf{IPFS} & \textbf{Arweave} & \textbf{传统 CDN} \\
\midrule
核心激励 & 检索 (payment) & 存储 (PoSt) & 无 & 一次性存储费 & 服务合同 \\
\addlinespace
复制模式 & 自组织 (市场) & 协议强制 & 无 (利他) & 基金资助 & 中心化管理 \\
\addlinespace
存储证明 & ECDH+Merkle (ms) & zk-SNARK (GPU 数小时) & 无 & SPoRA & 无 (信任) \\
\addlinespace
原生货币 & MNT (运营商) + BSV (用户) & FIL & 无 & AR & 法币 \\
\addlinespace
用户支付 & BSV (payment 微支付) & FIL & 免费 & 上传后免费 & 订阅/带宽 \\
\addlinespace
区块链 & BSV 同构侧链 & 自定义 (Lotus) & 无 & 自定义 & 无 \\
\addlinespace
内容加密 & 默认 (Method~42) & 可选 & 无 & 无 & 仅 TLS \\
\addlinespace
数据商务 & HTLC 原子交换 & 无 & 无 & 无 & DRM / 付费墙 \\
\addlinespace
地理分布 & 市场驱动 (利润) & 无激励 & 无激励 & 无激励 & 中心化优化 \\
\addlinespace
冷数据策略 & 一对一存档合约 & 扇区交易 & Pin 服务 & 永久存储 & 源站服务器 \\
\bottomrule
\end{tabularx}
\caption{去中心化存储和分发系统对比}
\end{table}

根本区别在于激励的指向。Filecoin 付费让节点证明它们存储了数据。Metanet 网络付费让节点服务数据。这种激励导向的差异产生不同的网络行为：Filecoin 为数据持久性优化，而 Metanet 网络为数据可用性和分发性能优化。

另外两个差异化要点值得关注。第一，BSV 同构链设计意味着 Metanet 网络不引入新的密码学原语、自定义虚拟机或新的脚本语言。开发者和节点运营商使用他们已经熟悉的相同 Bitcoin 交易格式和 Script 引擎。第二，双币种模型意味着终端用户完全与 CDN 的内部经济机制隔离。下载文件的用户支付 BSV 并接收数据。他们不需要知道 Metanet Chain 的存在、MNT 代币在运营商之间流通、或者存档合约正在侧链上结算。这种分离对主流采用至关重要。

% -----------------------------------------------------------------------
\section{11. metanet CLI}

运营 Metanet 节点通过 \texttt{metanet} 命令行界面管理。CLI 提供节点生命周期管理、合约管理和可选挖矿的最小命令集。

\begin{lstlisting}
metanet init          初始化新的 Metanet 节点 (生成密钥, 配置存储)
metanet start         启动节点 daemon (开始服务内容并参与网络)
metanet stop          优雅关闭节点
metanet status        显示节点状态: 运行时间, 对等节点, 缓存文件, 收入, 合约

metanet contracts     列出活跃存档合约, 存储使用情况, 证明计划
metanet peers         显示已连接的对等节点, 网络拓扑, 带宽统计
metanet mine          启用或配置合并挖矿 (需要连接到 BTC/BSV 矿池)
\end{lstlisting}

Metanet 节点本质上是运行在提供者模式下的 BitFS daemon。它通过 HTTP 提供加密内容，处理 payment 支付，维护存档合约，提交存储证明，并可选地参与合并挖矿。\texttt{metanet} CLI 用网络特定的配置和监控功能封装了这些能力。

节点运营商做两个关键决策：缓存什么内容（基于预期的检索收入）以及接受多少存档合约（基于可用的存储容量和 MNT 定价）。两个决策都是经济性的，CLI 提供做出这些决策所需的指标：每文件收入、存储成本、带宽利用率和合约盈利能力。

运行 Metanet 节点的准入门槛被刻意设计得很低。任何拥有存储容量、带宽和 BSV 钱包的机器都可以参与。没有最低质押要求，没有罚没机制，没有复杂的设置流程。节点运营商下载软件，初始化节点，然后开始缓存他们预期有利可图的内容。随着更多运营商发现服务热门内容是一项可行的生意，网络有机增长。

% -----------------------------------------------------------------------
\section{12. 结论}

Metanet 网络弥补了去中心化存储生态系统中的一个空白，将经济激励与用户真正看重的行为对齐：快速、可靠地访问数据。

核心机制是基于检索的激励。通过允许 Metanet 节点从服务内容中赚取 payment 费用，网络创建了一个正反馈循环，热门内容被广泛复制且随时可访问，无需任何中心化协调或协议管理的复制。这种自组织的 CDN 行为自然地从各个节点追求理性自利中涌现。

对于无法通过检索费维持自身的内容，存档合约提供了一种简单的双边机制来保证存储。使用 ECDH 重加密进行存储证明——而非零知识证明——将计算开销降低了数个数量级，同时保留了每份副本密码学唯一且可验证的基本属性。

双币种模型保持了系统的可访问性。用户以 BSV 支付，这是他们已经持有的货币。MNT 代币在网络运营商之间流通，其价值反映 CDN 服务的需求。这种分离确保 CDN 的内部经济机制不会给终端用户造成摩擦。

与提供文件系统协议、加密和所有权语义的 BitFS 一起，Metanet 网络完成了一个去中心化内容平台：数据由其创建者拥有，默认加密，可通过去信任的原子交换购买，由经济驱动的节点网络分发。整个链条的任何环节都不需要可信中介。

系统的设计在技术选择上刻意保守。BSV 同构侧链而非新型区块链。Bitcoin Script 而非自定义虚拟机。ECDH 和 Merkle 证明而非零知识密码学。合并挖矿而非新的共识机制。每一个选择都降低复杂性、构建在经过验证的基础设施之上、并降低参与门槛。贡献不在于组件本身，而在于它们如何被组装：一个协议奖励正确行为的 CDN，网络从个体经济决策中涌现，而非中心化规划。

% -----------------------------------------------------------------------
\section*{参考文献}

\begin{enumerate}
\item BitFS 项目, ``BitFS: 区块链上的点对点加密文件系统,'' 2025.

\item nChain, ``The Metanet Technical Summary v1.0,'' 2020. 基于区块链的有向无环图，用于组织数据关系.

\item C.~S.~Wright, ``An Immutable File and Data Store,'' nChain, 2019. Method~42: 使用 ECDSA/secp256k1 的确定性密钥派生实现逐文件加密.

\item S.~Nakamoto, ``Bitcoin: A Peer-to-Peer Electronic Cash System,'' 2008.

\item GB2608179A, ``Multi-level Blockchain,'' UK Intellectual Property Office, 2021. 在核心区块链上嵌入二级数据链.

\item P.~Wuille, ``BIP32: Hierarchical Deterministic Wallets,'' Bitcoin Improvement Proposals, 2012.

\item Protocol Labs, ``Filecoin: A Decentralized Storage Network,'' 2017.
\end{enumerate}

\end{document}
